\section{Iterative Postorder Traversal (One Stack)}
\label{sec:trees-iter-post-one}

\problemstatement{Return the postorder traversal of a binary tree without
recursion, using a single stack and $O(h)$ auxiliary space.}

\begin{patternbox}
The tell-tale that a one-stack postorder is wanted is a space constraint:
the two-stack method is simpler but stores all $n$ nodes. With one stack we
must handle the genuine difficulty head-on -- \emph{a node may sit on the
stack while we visit first its left subtree and only later its right} --
so we need a way to know, when a node resurfaces, whether we are returning
from its left or its right side.
\end{patternbox}

\subsection*{Understanding the problem carefully}

In postorder we record a node only after \emph{both} its children are
done. When we pop back up to a node, we must distinguish ``I have finished
its left subtree, now go right'' from ``I have finished its right subtree,
now record it.'' Tracking the \textbf{last node we recorded} resolves this
ambiguity: if that last-recorded node is the current node's right child
(or the right child is null), both subtrees are complete and the node may
be recorded; otherwise there is still a right subtree to descend into.

\begin{ideabox}[Remember the Last Node You Output]
Keep \textit{curr} (the cursor) and \textit{lastVisited} (the most
recently recorded node). Descend left, pushing. At the top of the stack:
if its right child exists and has not just been visited, move right;
otherwise record the top, pop it, and set \textit{lastVisited} to it. That
single comparison against \textit{lastVisited} tells you which side you are
returning from.
\end{ideabox}

\subsection*{Algorithm}

\begin{enumerate}
  \item \textit{curr} $=$ root; \textit{lastVisited} $=$ null; empty
        stack.
  \item While \textit{curr} is non-null or the stack is non-empty:
  \begin{enumerate}
    \item While \textit{curr} is non-null: push \textit{curr},
          \textit{curr} $=$ \textit{curr}.left.
    \item Let \textit{peek} $=$ top of stack. If \textit{peek}.right
          exists and \textit{peek}.right $\ne$ \textit{lastVisited}: set
          \textit{curr} $=$ \textit{peek}.right.
    \item Else: record \textit{peek}, pop it, set \textit{lastVisited}
          $=$ \textit{peek}.
  \end{enumerate}
\end{enumerate}

\subsection*{Dry run}

\dryrun{Tree: root $1$; left $2$; right $3$.}

\begin{itemize}
  \item Dive left from $1$: push $1$, push $2$. \textit{curr}$=$null.
  \item Peek $2$: right is null $\Rightarrow$ record $2$, pop,
        \textit{lastVisited}$=2$.
  \item Peek $1$: right $3\ne$\textit{lastVisited} $\Rightarrow$
        \textit{curr}$=3$. Push $3$.
  \item Peek $3$: right null $\Rightarrow$ record $3$, pop,
        \textit{lastVisited}$=3$.
  \item Peek $1$: right $3=$\textit{lastVisited} $\Rightarrow$ record $1$,
        pop.
  \item Result: $2,3,1$.
\end{itemize}

\subsection*{C++ implementation}

\begin{lstlisting}
#include <bits/stdc++.h>
using namespace std;

struct TreeNode {
    int val;
    TreeNode* left;
    TreeNode* right;
    TreeNode(int v) : val(v), left(nullptr), right(nullptr) {}
};

vector<int> postorderOneStack(TreeNode* root) {
    vector<int> out;
    stack<TreeNode*> st;
    TreeNode* curr = root;
    TreeNode* lastVisited = nullptr;

    while (curr != nullptr || !st.empty()) {
        while (curr != nullptr) {
            st.push(curr);
            curr = curr->left;
        }
        TreeNode* peek = st.top();
        if (peek->right != nullptr && peek->right != lastVisited) {
            curr = peek->right;             // still owe the right subtree
        } else {
            out.push_back(peek->val);       // both sides done: record
            lastVisited = peek;
            st.pop();
        }
    }
    return out;
}

int main() {
    TreeNode* root = new TreeNode(1);
    root->left = new TreeNode(2);
    root->right = new TreeNode(3);

    for (int x : postorderOneStack(root)) cout << x << " ";
    cout << "\n";
    return 0;
}
\end{lstlisting}

\begin{complexitybox}
\textbf{Time Complexity.} $O(n)$ -- each node is pushed once, and the
peek/right check amortises to $O(1)$ per node. \textbf{Space Complexity.}
$O(h)$ -- the stack only ever holds the current root-to-node path, unlike
the two-stack version's $O(n)$.
\end{complexitybox}

\begin{takeawaybox}
When you must return from a node's left versus right side and act
differently, a \textit{lastVisited} pointer is the standard way to tell
them apart without extra per-node flags. This ``remember where you came
from'' idea generalises to any iterative post-processing traversal.
\end{takeawaybox}
